\section{Introduction}
A wiring diagram constrains which neurons can influence each other, but it does not specify a complete behaving animal. A computational experiment must still choose how external events become neural inputs, which dynamical parameters to use, what state persists between trials, and how activity is interpreted. These choices become especially important when a model is presented through an interactive language interface: a fluent explanation can hide whether a statement follows from neural activity, a supplied behavioral summary, or an assumption in the interface.

We assemble these choices into an interactive research platform for a connectome-based fly model. A user can specify a stimulus, inspect the neural and motor outputs, alter the model's experience, and compare subsequent responses. The system joins an adult-brain spiking simulation to measured odor-response profiles, additional sensory channels, persistent mushroom-body synaptic factors, a cross-connectome motor-output bridge, and a language interface. Its scientific object is the assembled model and the experiments that can be reproduced with it.

The contribution has three parts: an integrated experimental interface with explicit component provenance; three recorded case studies connecting interventions to neural and motor outputs; and a characterization of recorded language interactions and the auxiliary neural-state reader. The cases establish repeatable chemical representations, distinct sensory-output patterns, and recoverable changes in odor responses after experience. They also expose a meaningful boundary: separating external odor and reinforcement pulses does not prevent learning updates when odor input itself recruits dopamine neurons. The system makes such dependencies available for inspection and further experiments. Its claims concern the assembled model and tested conditions, rather than a new learning algorithm or a validated whole animal.
